% !TeX root = main.tex

\section{实验}
\subsection{数据集}

我们在三个自然语言被动聆听数据集上进行实验，包括一个 EEG 数据集和两个 MEG 数据集。
三个数据集覆盖中文和法语两种语言，并包含连续故事聆听和多故事自然语言聆听两类实验范式。

ChineseEEG-2 包含自然中文阅读与聆听条件。
本研究使用其中 Passive Listening 条件下的 8 名受试者，
并仅分析《小王子》刺激材料。
受试者聆听由真实朗读产生的中文语音，
脑活动由 128 通道 EEG 系统记录。
由于实际朗读可能包含相对于原始文本的停顿、漏读、增读或替换，
我们进一步核对实际播放音频，并基于修正后的文本获得词级时间戳。

SMN4Lang 是一个面向自然语言理解研究的中文多模态神经数据集。
完整数据包含 12 名中文母语受试者聆听 60 个自然故事时的 MEG 和 fMRI 记录。
每个故事持续约 4--7~min，内容覆盖多个不同主题。
本研究使用其中的 MEG 数据，并以故事中提供的词级时间对齐信息构建词事件。

LittlePrince\_MEG\_French\_Listen\_Pallier2025
包含法语母语受试者连续聆听法语有声书《小王子》时记录的 MEG 数据。
完整故事被组织为 9 个连续 listening runs。
MEG 由 306 个传感器记录，包括 102 个 magnetometers
和 204 个 planar gradiometers。
本研究当前实验使用其中的 sub-01--sub-10。

对于数据集作者已经提供预处理衍生数据的情况，
我们直接使用官方预处理后的神经信号；
对于其余数据，仅进行与原研究一致的基础处理以及统一输入格式所需的最小转换。
为了保持三个数据集后续词级解码输入的一致性，
所有神经信号最终统一至 50~Hz 的时间采样率。

\subsection{词级神经分词}

我们将三个数据集统一转换为词级神经解码任务。
对于每一个词事件，以其语音起始时间（word onset）作为时间零点，
截取 onset 后连续 1~s 的神经活动作为当前词的模型输入。
因此，在统一的 50~Hz 采样率下，
每一个词级样本包含 50 个时间点。

记第 $i$ 个词的起始时间为 $t_i$，
连续神经记录为 $\mathbf{X}(t)$，
则该词对应的神经输入定义为

\[
\mathbf{X}_i =
\mathbf{X}[t_i,\;t_i + 1~\mathrm{s}).
\]

我们在三个数据集上均采用相同的 1~s 输入窗口，
以避免由于神经窗口长度不同而引入额外的跨数据集差异。

对于 ChineseEEG-2，
词 onset 根据实际播放的朗读音频确定，
而非原始视觉呈现时间。
我们首先核对实际朗读内容与原始文本之间的差异，
随后通过强制对齐获得字符级和词级时间戳，
并据此提取相应的 EEG 窗口。

SMN4Lang 和 Pallier2025 则利用数据集中已有的词级时间信息，
将其转换至对应 MEG recording 的时间坐标后构建词级样本。


\subsection{文本表示提取}

我们将词汇解码建模为神经表示与预训练文本表示之间的对齐问题，
而不是直接对离散词类别进行分类。
对于每一个候选词，
我们首先使用冻结的预训练文本模型提取其词表示，
随后训练神经编码器将对应的 EEG/MEG 窗口映射到相同的表示空间。

对于两个中文数据集 ChineseEEG-2 和 SMN4Lang，
我们使用 Mengzi-T5-base 作为文本编码器。
对于法语 Pallier2025 数据集，
我们使用 T5-large。
两种模型在神经解码训练过程中均保持冻结。

设文本编码模型共有 $L$ 个隐藏层。
我们提取模型中间深度位置的隐藏状态，
即约 $0.5L$ 处的表示。
当一个词被分解为多个 subword tokens 时，
我们对所有有效 token 的隐藏状态进行平均，
得到该词的固定维度表示：

\[
\mathbf{y}_i =
\frac{1}{K_i}
\sum_{k=1}^{K_i}
\mathbf{h}_{i,k},
\]

其中 $K_i$ 表示第 $i$ 个词包含的有效 token 数量，
$\mathbf{h}_{i,k}$ 表示对应 token 在所选隐藏层的表示。

ChineseEEG-2 和 SMN4Lang 使用 768 维文本表示，
Pallier2025 使用 1024 维文本表示。
所有词表示在模型训练前预先提取并冻结，
神经编码器的目标是预测与真实词相对应的文本向量。


\subsection{数据划分}

三个数据集均采用未见刺激划分，
即训练集、验证集和测试集按照连续语言材料进行分离，
而不是随机拆分单个词或句子。
对于同一数据集，材料级划分在所有受试者之间保持一致，
因此某一语言材料一旦被分配到测试集，
就不会以其他受试者的记录形式再次出现在训练集中。

对于 ChineseEEG-2，
我们以《小王子》的章节作为划分单位。
第 1、4、6--8、10--13、15--24 和 26 章用于训练，
第 3、5 和 25 章用于验证，
第 2 和 9 章保留为测试。
由于其中一名受试者的记录存在截断，
第 14 章和第 27 章在所有受试者中统一排除，
以保证多受试者之间使用相同的有效材料范围。

对于 SMN4Lang，
我们以完整故事作为划分单位。
第 1--50 个故事用于训练，
第 51--55 个故事用于验证，
第 56--60 个故事保留为测试。
这样可以保证测试阶段使用的故事内容
在模型训练和模型选择过程中均未出现。


\subsection{词表构建}

本研究将模型训练所使用的词事件与最终评价使用的候选词表分开定义。
在训练阶段，我们使用训练集中所有满足词级窗口要求且具有有效文本表示的词事件进行对比学习，
而不将模型限制为固定数量的词类别。
因此，模型学习的是神经活动与连续文本表示空间之间的映射，
而不是一个固定词表上的多分类器。

为了与已有词汇解码研究中的闭集评价保持可比性，
我们首先在每个数据集的训练材料中构建高频 50 词候选词表，
并将其作为主要的闭集评价基线。
候选词按照训练材料中的出现频率从高到低排序；
当多个词具有相同频次时，
使用固定的词典序规则确定其顺序，
以保证词表构建过程具有确定性。

由于三个数据集均包含多个受试者聆听相同语言材料，
词频统计以语言材料中的词出现为单位，
而不是直接累加所有受试者的神经事件。
换言之，
同一刺激中的一个词即使被多名受试者重复聆听，
在构建候选词表时也只按照其在语言材料中的出现次数计数，
从而避免受试者数量人为改变词频分布。

对于 ChineseEEG-2，
男性和女性实际朗读版本分别被视为独立的刺激序列。
由于两个版本在实际朗读过程中可能存在漏读、增读或替换，
词频统计基于各自经过核对后的实际朗读内容，
但同一朗读版本被多名受试者重复聆听不会重复增加词频。
SMN4Lang 和 Pallier2025 同样按照语言材料本身的词事件统计词频，
而不按照受试者重复次数计数。

固定的 50 词词表能够提供直观的闭集检索基线，
但较小的候选词表通常更容易获得较高的 Top-$k$ 准确率，
同时只能覆盖真实语言中的有限部分。
因此，
我们进一步从相同的训练材料中构建
$N\in\{20,50,100,150\}$ 的多尺度候选词表。
不同规模词表均遵循相同的训练集词频排序规则，
且在测试集评价之前冻结。

测试集中的词频、词汇覆盖率和解码结果均不参与候选词的选择。
这一设计使我们能够在保持候选词来源一致的条件下，
进一步分析词表规模、闭集检索性能与实际可传递语言信息之间的关系。


\subsection{评价指标}

我们首先使用 Top-$1$ 和 Top-$10$ 检索准确率评价模型在固定候选词表中的词汇辨识能力。
对于每一个词级神经样本，
模型输出一个神经表示，
并计算其与候选词文本表示之间的余弦相似度。
随后按照相似度从高到低对候选词进行排序。

记第 $i$ 个样本对应的真实词为 $y_i$，
模型给出的候选词排序为
$\hat{y}_{i,1},\hat{y}_{i,2},\ldots$，
则 Top-$k$ 准确率定义为

\[
\mathrm{Top}\text{-}k
=
\frac{1}{N}
\sum_{i=1}^{N}
\mathbb{I}
\left(
y_i
\in
\{\hat{y}_{i,1},\ldots,\hat{y}_{i,k}\}
\right),
\]

其中 $N$ 表示评价样本数量，
$\mathbb{I}(\cdot)$ 为指示函数。
本研究主要报告 Top-$1$ 和 Top-$10$。

由于自然语言中的词频分布高度不均衡，
直接对所有词事件求平均会使高频词对最终指标产生更大影响。
因此，
除事件级微平均准确率外，
我们进一步计算词类型宏平均准确率：
首先分别计算每一个真实词的 Top-$k$ 命中率，
随后对不同词类型赋予相同权重进行平均。
主模型选择使用固定 50 词候选词表上的宏平均 Top-$10$ 准确率，
以减少高频功能词对模型选择的支配作用。

然而，
Top-$k$ 准确率只描述模型在一个预先给定的候选词表内部能够多好地区分词，
并不能反映该词表能够覆盖真实语言的多少。
例如，
一个只包含少量高频词的系统可能具有很高的 Top-$10$ 准确率，
但仍无法表达故事中的大部分词汇。
因此，
我们进一步使用开放词表互信息
（Open-Vocabulary Mutual Information, OVMI）
衡量解码器实际能够传递的语言信息。

设 $\Omega$ 为参考语言分布包含的完整词空间，
$p(x)$ 表示用户在目标语言环境中产生词 $x$ 的概率，
$S\subseteq\Omega$ 为解码器支持的候选词表。
词表覆盖率定义为

\[
C(S)
=
\Pr(X\in S)
=
\sum_{x\in S} p(x).
\]

在候选词表内部，
模型真实词 $X$ 与预测词 $Y$ 之间的互信息记为

\[
I(X;Y\mid X\in S).
\]

开放词表互信息定义为

\[
I_{\mathrm{OVMI}}(S)
=
C(S)
I(X;Y\mid X\in S).
\]

因此，
该指标同时考虑两个互补因素：
候选词表能够覆盖目标语言分布的比例，
以及模型在已支持词汇内部实际能够传递的信息。
当候选词表扩大时，
词表内的检索任务通常会变得更困难，
但语言覆盖率也会同时增加；
开放词表互信息使这两种变化能够在统一的信息尺度上进行比较。

本研究分别使用故事内参考分布和领域参考分布计算开放词表互信息。
故事内参考由当前目标语言材料自身的词频分布构建，
用于回答模型在特定故事场景下能够传递多少语言信息。
对于 ChineseEEG-2 和 Pallier2025，
该参考对应《小王子》的目标文本分布；
对于 SMN4Lang，
则由其完整自然故事材料构建相应参考分布。

领域参考使用独立于当前实验材料的外部语料构建，
用于回答相同解码能力在更广泛目标语言环境中能够覆盖多少信息。
故事内参考强调针对特定叙事内容的通信能力，
而领域参考更强调词汇能力向目标应用场景的外推。
两种参考分布均只参与最终评价，
不参与模型训练和候选词表构建。

开放词表互信息以 bit 为单位报告。
其计算基于固定候选词表上的 Top-$1$ 预测所形成的混淆矩阵，
并结合相应参考语言分布估计词表覆盖率和词表内互信息。
当某一候选词在评价集中缺少真实样本，
因而无法可靠估计对应混淆矩阵行时，
我们不通过删除候选词或人工平滑来补全该信息，
而将完整开放词表互信息标记为不可计算。



\subsection{严格审计}

较高的词汇检索性能并不必然意味着模型利用了当前词对应的神经活动。
在连续语言任务中，
同一 recording 内的时间位置、上下文长度、材料顺序以及相邻词之间的结构关系，
都可能为目标词提供额外的预测线索。
这一问题在引入句子级上下文后更加突出：
即使当前神经输入与目标词不再正确对应，
模型仍可能依赖上下文结构缩小候选范围。
因此，
除标准解码条件外，
我们进一步构建一组受控审计实验，
以区分正确神经输入带来的增量与非神经结构先验所支持的性能。

标准条件保留每个目标词与其真实神经窗口之间的对应关系，
并作为所有审计条件的参照。
对于所有对照实验，
我们尽可能保持模型参数、候选词表、目标词集合以及上下文位置不变，
仅改变需要检验的信息来源。
因此，
标准条件与对照条件之间的差异可以用于衡量被破坏信息对模型性能的贡献。

首先，
我们使用时间错位对照破坏目标词与神经活动之间的精确时间对应。
对于每个目标词，
从同一受试者、同一 recording 中选取与目标窗口时间上不重叠的神经片段作为输入，
同时保持目标词标签及其在上下文序列中的位置不变。
这一设计尽可能保留同一 recording 中较慢的信号状态、受试者特征和整体实验环境，
但移除当前词与当前神经活动之间的局部时间对应。
因此，
如果正确时间对齐条件相对于时间错位条件仍具有稳定优势，
则表明模型性能中存在依赖精确神经--语言时间对应的成分。

其次，
我们构建错误神经供体对照，
进一步检验模型是否利用了输入神经窗口中与词内容相关的信息。
对于每一个目标事件，
供体神经窗口取自同一受试者和同一 recording 中的另一个词事件。
供体窗口不得与目标窗口重叠，
且供体词必须与目标词不同；
目标词标签、候选词表以及上下文位置均保持不变。
为了降低单次随机配对带来的偶然性，
供体替换使用多个预先确定的随机种子重复进行，
并在各条件共同支持的查询集合上进行比较。

错误供体实验除了比较总体解码性能外，
还可以检验预测是否跟随供体神经内容发生系统性变化。
记目标词为 $y_i$，
供体窗口对应的真实词为 $d_i$。
在替换神经输入后，
我们统计供体词 $d_i$ 在模型预测排序中的变化，
以及其进入 Top-$k$ 候选的概率。
如果错误神经输入不仅降低目标词的排名，
同时提高供体词的排名，
则说明模型输出确实受到所输入神经内容的驱动，
而不仅仅是在移除正确输入后产生无方向的性能下降。

最后，
针对句子级上下文模型，
我们设置不包含有效当前神经内容的结构对照。
该条件保留序列长度、词在上下文中的相对位置以及上下文分组方式等结构信息，
但移除可用于识别当前词的神经信号。
这一对照用于估计仅依赖任务结构和上下文组织时，
模型能够获得多高的词汇检索性能。
因此，
标准上下文模型的总性能可以进一步区分为
结构信息已经能够解释的部分，
以及在恢复正确神经输入后额外获得的部分。

所有审计条件均使用与标准条件相同的评价指标，
包括 Top-$1$、Top-$10$、宏平均 Top-$10$
以及在条件允许时计算的开放词表互信息。
对于信息量指标，
我们将正确神经输入条件与预先规定的对照条件之间的差值写为

\[
\Delta I_{\mathrm{align}}
=
I_{\mathrm{clean}}
-
I_{\mathrm{control}},
\]

其中 $I_{\mathrm{clean}}$ 表示正确神经--词对应条件下的信息量，
$I_{\mathrm{control}}$ 表示相应审计条件下的信息量。
本文将该差值解释为
\emph{神经时间对齐特异的信息增量}，
即恢复正确神经--语言对应关系后增加的信息，
而不将其直接称为``纯神经信息''，
因为任何单一对照均不能排除所有可能的非神经因素。